\chapter{วิทยาศาสตร์การออกกำลังกายและชีวกลศาสตร์สุขภาพประยุกต์}
\label{chap:ch03}

\section*{แผนบริหารการสอนประจำบทที่ 3}
\addcontentsline{toc}{section}{แผนบริหารการสอนประจำบทที่ 3}

\noindent\textbf{หัวข้อเนื้อหาประจำบท}
\begin{enumerate}
    \item อัตราความถี่ก้าวและการจำแนกความหนักของการออกกำลังกายตามเกณฑ์สากล
    \item สรีรวิทยาการคำนวณระยะก้าวเดินแบบแปรผันตามเพศและความสูงของผู้ใช้
    \item การประเมินพลังงานเผาผลาญพื้นฐานด้วยสมการมาตรฐาน Mifflin-St Jeor
    \item การคำนวณพลังงานที่ใช้ทั้งหมดในแต่ละวันและการเผาผลาญจากกิจกรรมก้าวเดิน
    \item กลไกสรีรวิทยาของพฤติกรรมเนือยนิ่งและระบบแจ้งเตือนการหยุดนั่งนาน
    \item การประเมินสมดุลของเหลวในร่างกายและความต้องการน้ำบริสุทธิ์ในแต่ละวัน
\end{enumerate}

\noindent\textbf{วัตถุประสงค์เชิงพฤติกรรม}
\begin{enumerate}
    \item คำนวณและจำแนกระดับความหนักของกิจกรรมทางกายจากอัตราความถี่ก้าว (Cadence) ได้อย่างถูกต้อง
    \item ประยุกต์ใช้สมการสัดส่วนทางชีวกลศาสตร์ในการประเมินระยะก้าวและระยะทางรวมได้อย่างถูกต้อง
    \item คำนวณอัตราการเผาผลาญพื้นฐาน (BMR) และการใช้พลังงานรวม (TDEE) ตามสูตร Mifflin-St Jeor ได้อย่างแม่นยำ
    \item อธิบายผลเสียทางสรีรวิทยาของการนั่งนิ่งติดต่อกันเกิน 1 ชั่วโมง พร้อมแนวทางแก้ไขได้อย่างถูกต้อง
    \item คำนวณปริมาณน้ำดื่มที่เหมาะสมตามน้ำหนักตัวและระดับกิจกรรมทางกายได้อย่างถูกต้อง
\end{enumerate}

\noindent\textbf{กิจกรรมการเรียนการสอน}
\begin{enumerate}
    \item บรรยายหลักการเวชศาสตร์การออกกำลังกาย สรีรวิทยาการเผาผลาญ และชีวกลศาสตร์การเคลื่อนไหว
    \item กิจกรรมฝึกคำนวณ BMR, TDEE, Stride Length และ Hydration จากโจทย์กรณีศึกษาบุคคลตัวอย่าง
    \item การทดสอบเดินจับเวลาเพื่อวัด Cadence จริงเปรียบเทียบกับค่าที่แอปพลิเคชันคำนวณได้
    \item การอภิปรายกลุ่มและการตอบคำถามทบทวนท้ายบท
\end{enumerate}

\noindent\textbf{สื่อการเรียนการสอน}
\begin{enumerate}
    \item เอกสารคำสอนบทที่ 3 วิทยาศาสตร์การออกกำลังกายและชีวกลศาสตร์สุขภาพ
    \item เครื่องชั่งน้ำหนักและสายวัดความสูงมาตรฐาน
    \item นาฬิกาจับเวลาและสมาร์ทโฟนที่ติดตั้งแอปพลิเคชัน Smart Health Tracker
\end{enumerate}

\noindent\textbf{การประเมินผล}
\begin{enumerate}
    \item ประเมินความถูกต้องจากการทำใบงานคำนวณสรีรวิทยาการเผาผลาญ
    \item ประเมินทักษะการวัดและสอบเทียบระยะก้าวเดินในการปฏิบัติการภาคสนาม
    \item ประเมินผลสัมฤทธิ์จากคำถามทบทวนท้ายบทที่ 3
\end{enumerate}

\newpage

\section{อัตราความถี่ก้าวและการจำแนกความหนักของการออกกำลังกาย}
ในทางเวชศาสตร์การออกกำลังกาย อัตราความถี่ก้าว (Cadence) มีหน่วยเป็นก้าวต่อนาที (Steps per Minute: SPM) เป็นดัชนีชี้วัดที่สะท้อนถึงความหนักของการออกกำลังกาย (Exercise Intensity) ได้อย่างตรงไปตรงมา การนับจำนวนก้าวสะสมเพียงอย่างเดียวโดยไม่คำนึงถึง Cadence ไม่อาจบ่งชี้ได้ว่ากิจกรรมนั้นส่งเสริมสมรรถภาพของระบบหัวใจและหลอดเลือด (Cardiorespiratory Fitness) หรือไม่

\begin{healthbox}
\textbf{เกณฑ์การจำแนกความหนักของการก้าวเดินตามมาตรฐานวิทยาลัยเวชศาสตร์การกีฬาแห่งสหรัฐอเมริกา (ACSM)}

\begin{itemize}
    \item \textbf{การเดินช้าหรือพฤติกรรมเนือยนิ่ง ($<80\text{ SPM}$)} มีการเผาผลาญพลังงานในระดับเบามาก เหมาะสำหรับผู้ฟื้นฟูสมรรถภาพหรือผู้สูงอายุ
    \item \textbf{การเดินเพื่อสุขภาพระดับปานกลาง ($100 - 120\text{ SPM}$)} เป็นช่วงความถี่ที่เทียบเท่ากับความหนักระดับปานกลาง (Moderate Intensity, $3.0 - 5.9\text{ METs}$) ซึ่งทำให้หัวใจเต้นในย่านแอโรบิก ช่วยลดระดับน้ำตาลในเลือดและความดันโลหิต
    \item \textbf{การเดินเร็วหรือการวิ่งเหยาะ ($>130\text{ SPM}$)} จัดเป็นกิจกรรมทางกายระดับหนัก (Vigorous Intensity, $\ge 6.0\text{ METs}$) ช่วยเพิ่มความจุปอดและเสริมสร้างความแข็งแรงของกล้ามเนื้อหัวใจ
\end{itemize}
\end{healthbox}

\section{ชีวกลศาสตร์การคำนวณระยะก้าวเดินเฉพาะบุคคล}
แอปพลิเคชันนับก้าวทั่วไปมักกำหนดระยะก้าวเดินเป็นค่าคงที่เฉลี่ย เช่น $70\text{ ซม.}$ ส่งผลให้ผู้ใช้งานที่มีความสูงต่างกันได้รับค่าระยะทางรวมที่ผิดเพี้ยนอย่างมีนัยสำคัญ

ในทางมานุษยมิติ (Anthropometry) ความยาวของกระดูกโคนขาและขาท่อนล่างมีความสัมพันธ์เชิงเส้นตรงกับความสูงของร่างกาย ($H$) นอกจากนี้สรีระของเพศชายและเพศหญิงยังมีมุมของกระดูกเชิงกราน (Q-angle) และจุดหมุนของข้อสะโพกที่แตกต่างกัน ระบบ Smart Health Tracker จึงเลือกใช้แบบจำลองสัดส่วนตามเพศและความสูงที่ผ่านการยอมรับในวารสารชีวกลศาสตร์สากล

\begin{equation}
L_{\text{stride}} = 
\begin{cases}
H \times 0.415 & \text{สำหรับเพศชาย} \\
H \times 0.413 & \text{สำหรับเพศหญิง}
\end{cases}
\end{equation}
เมื่อ $H$ คือความสูงของผู้ใช้ในหน่วยเซนติเมตร และ $L_{\text{stride}}$ คือความยาวก้าวในหน่วยเซนติเมตร ระยะทางรวมสะสม ($D$) ในหน่วยกิโลเมตรจึงคำนวณได้จาก
\begin{equation}
D = \frac{\text{Steps} \times L_{\text{stride}}}{100{,}000}
\end{equation}

\begin{table}[H]
\centering
\caption{ตัวอย่างระยะก้าวเดินจำเพาะบุคคลตามความสูงและเพศสภาพ}
\label{tab:personalized_stride}
\small
\begin{tabularx}{\textwidth}{|c|c|c|X|}
\toprule
\textbf{ความสูง ($H$)} & \textbf{ระยะก้าวเพศชาย ($0.415$)} & \textbf{ระยะก้าวเพศหญิง ($0.413$)} & \textbf{ระยะทางสะสมเมื่อเดิน 10,000 ก้าว} \\
\midrule
$150\text{ cm}$ & $62.25\text{ cm}$ & $61.95\text{ cm}$ & ชาย $6.23\text{ km}$ \textbar\ หญิง $6.20\text{ km}$ \\
$160\text{ cm}$ & $66.40\text{ cm}$ & $66.08\text{ cm}$ & ชาย $6.64\text{ km}$ \textbar\ หญิง $6.61\text{ km}$ \\
$170\text{ cm}$ & $70.55\text{ cm}$ & $70.21\text{ cm}$ & ชาย $7.06\text{ km}$ \textbar\ หญิง $7.02\text{ km}$ \\
$180\text{ cm}$ & $74.70\text{ cm}$ & $74.34\text{ cm}$ & ชาย $7.47\text{ km}$ \textbar\ หญิง $7.43\text{ km}$ \\
\bottomrule
\end{tabularx}
\end{table}

\section{การประเมินการเผาผลาญพลังงานด้วยสมการ Mifflin-St Jeor}
พลังงานที่ร่างกายมนุษย์ใช้ในแต่ละวันประกอบด้วยพลังงานที่ใช้ในการทำงานของอวัยวะภายในขณะพักผ่อน (Basal Metabolic Rate: BMR) รวมกับพลังงานที่ใช้ในกิจกรรมทางกายและการย่อยอาหาร (Total Daily Energy Expenditure: TDEE)

ในบรรดาสมการประเมิน BMR สมาคมโภชนาการและการกำหนดอาหารแห่งสหรัฐอเมริกา (Academy of Nutrition and Dietetics) ได้รับรองว่าสมการของ Mifflin-St Jeor (1990) เป็นสมการที่มีความเที่ยงตรงและแม่นยำสูงสุดสำหรับประชากรทั่วไป

\begin{theorybox}
\textbf{สมการประเมินอัตราการเผาผลาญพลังงานพื้นฐานของ Mifflin-St Jeor}

กำหนดให้ $W$ คือน้ำหนักตัว (กิโลกรัม), $H$ คือความสูง (เซนติเมตร) และ $A$ คืออายุ (ปี)
\begin{align}
\text{BMR}_{\text{male}} &= 10W + 6.25H - 5A + 5 \\
\text{BMR}_{\text{female}} &= 10W + 6.25H - 5A - 161
\end{align}
การคำนวณพลังงานที่ใช้รวมตลอดวัน (TDEE) คำนวณได้จากการคูณ BMR ด้วยตัวคูณระดับกิจกรรมทางกาย (Physical Activity Level: PAL)
\begin{equation}
\text{TDEE} = \text{BMR} \times \text{PAL}
\end{equation}
โดยที่ $\text{PAL}$ สำหรับผู้มีกิจกรรมน้อย (Sedentary) มีค่าเท่ากับ $1.2$ กิจกรรมปานกลาง $1.55$ และกิจกรรมหนัก $1.725$
\end{theorybox}

สำหรับการเผาผลาญพลังงานเฉพาะเจาะจงจากกิจกรรมการก้าวเดิน (Exercise Activity Thermogenesis: EAT) ระบบคำนวณจากค่าสมมูลเมแทบอลิก (Metabolic Equivalent of Task: MET) ของการเดินที่ระดับความเร็วต่างๆ
\begin{equation}
\text{Calories (kcal)} = \text{Time (hr)} \times \text{MET} \times W\text{ (kg)}
\end{equation}
โดยการเดินที่ความเร็วปานกลาง ($100\text{ SPM}$) มีค่าประมาณ $3.5\text{ METs}$

\section{กลไกสรีรวิทยาของพฤติกรรมเนือยนิ่งและระบบแจ้งเตือน}
การนั่งนิ่งติดต่อกันเป็นเวลานานเกินกว่า 60 นาที ส่งผลให้การหดตัวของกล้ามเนื้อโครงร่างขนาดใหญ่โดยเฉพาะกล้ามเนื้อขาและสะโพกหยุดชะงักลง ก่อให้เกิดผลกระทบทางชีวเคมีร้ายแรง ได้แก่
\begin{enumerate}
    \item เอนไซม์ไลโพโปรตีนไลเปส (Lipoprotein Lipase: LPL) บริเวณกล้ามเนื้อหยุดทำงาน ส่งผลให้ความสามารถในการกำจัดไตรกลีเซอไรด์และไขมันในกระแสเลือดลดลงมากกว่าร้อยละ 90
    \item เกิดภาวะดื้อต่ออินซูลินแบบเฉียบพลัน (Acute Insulin Resistance) ทำให้เซลล์ไม่สามารถดึงน้ำตาลกลูโคสไปใช้ได้ตามปกติ
    \item การไหลเวียนของเลือดดำบริเวณขาลดลง ก่อให้เกิดความดันในหลอดเลือดและเสี่ยงต่อภาวะลิ่มเลือดอุดตันในหลอดเลือดดำส่วนลึก (Deep Vein Thrombosis: DVT)
\end{enumerate}

ระบบ Smart Health Tracker จึงได้ออกแบบระบบตรวจสอบพฤติกรรมเนือยนิ่ง (Sedentary Behaviour Engine) เมื่อเซนเซอร์ตรวจพบว่าไม่มีก้าวเดินเกิดขึ้นติดต่อกันครบ 60 นาที ระบบจะกระตุ้นการแจ้งเตือนทั้งในรูปแบบภาพสั่นไหวและข้อความเชิงบวก เพื่อให้ผู้ใช้ลุกขึ้นยืดเหยียดกล้ามเนื้อหรือเดินเปลี่ยนอิริยาบถอย่างน้อย $2 - 3$ นาที

\section{สมดุลของเหลวในร่างกายและการดื่มน้ำบริสุทธิ์}
น้ำคิดเป็นร้อยละ 60 ของน้ำหนักตัวในผู้ใหญ่ การขาดน้ำเพียงร้อยละ 1 ถึง 2 ของน้ำหนักตัว จะส่งผลให้ปริมาตรพลาสมาในเลือดลดลง ประสิทธิภาพการระบายความร้อนของร่างกายแย่ลง และความดันโลหิตตก

ระบบจึงคำนวณเป้าหมายการดื่มน้ำพื้นฐานประจำวัน (Baseline Fluid Requirement) ตามมาตรฐานการแพทย์
\begin{equation}
V_{\text{water, base}} = W\text{ (kg)} \times 32.5\text{ mL}
\end{equation}
และเพิ่มปริมาณน้ำชดเชยตามเหงื่อที่สูญเสียไปจากการก้าวเดินในอัตรา $350\text{ mL}$ ต่อการเดินสะสมทุกๆ $3{,}000$ ก้าว พร้อมแถบกราฟิกบันทึกการดื่มน้ำที่อัปเดตแบบเรียลไทม์

\section*{คำถามทบทวนท้ายบทที่ 3}
\addcontentsline{toc}{section}{คำถามทบทวนท้ายบทที่ 3}
\begin{enumerate}
    \item เหตุใดอัตราความถี่ก้าว (Cadence) จึงมีความสำคัญมากกว่าการนับจำนวนก้าวสะสมเพียงอย่างเดียว
    \item ชายอายุ 35 ปี สูง 175 ซม. หนัก 70 กก. จงคำนวณระยะก้าวเดิน ความยาวทางสะสมเมื่อเดิน 8,000 ก้าว และค่า BMR ตามสูตร Mifflin-St Jeor
    \item จงอธิบายการเปลี่ยนแปลงทางชีวเคมีและผลเสียต่อหลอดเลือดเมื่อร่างกายอยู่ในภาวะเนือยนิ่งติดต่อกันเกิน 60 นาที
    \item เพราะเหตุใดระบบจึงต้องคำนวณปริมาณน้ำดื่มชดเชยเพิ่มขึ้นตามจำนวนก้าวเดินสะสม
\end{enumerate}
\clearpage
